\section{Motion Tracking}
While motion detection identifies the presence of movement, \textbf{motion tracking} aims to 
follow specific objects across consecutive frames. 
This process involves establishing temporal correspondence by associating detected entities in the current 
frame with their instances in subsequent frames.

Motion tracking is essential across various domains, including:
\begin{itemize}
    \item \textbf{People monitoring:} analyzing crowd dynamics or tracking a specific individual;
    \item \textbf{Vehicle tracking:} monitoring the movement of vehicles in traffic surveillance;
    \item \textbf{Biological applications:} tracking the movement of animals or cells in scientific research;
\end{itemize}

The primary objective is often \textbf{object tracking}, which maintains the persistent identity of a target over time. 
Depending on the application, tracking can be performed in different dimensionalities:
\begin{itemize}
\item Tracking 2D coordinates in the image plane (e.g., using bounding boxes or centroids);
\item Tracking 3D coordinates in real-world space (requires multiple cameras or depth sensors);
\item Determining the pose or articulation of complex objects, such as human skeletal tracking.
\end{itemize}

\subsection{2D Tracking Methodologies}
In scenes with a single moving object, tracking can often be achieved by simply following the object's centroid. 
However, multi-object scenarios require sophisticated association algorithms to prevent identity switches. 
Common 2D tracking approaches include:
\begin{itemize}
    \item \textbf{Region-based:} Associating areas based on global features like color or texture.
    \item \textbf{Contour-based:} Following the evolving boundary or "snake" of an object.
    \item \textbf{Feature-based:} Tracking distinct interest points, such as Harris corners or SIFT features.
    \item \textbf{Template-based:} Matching a predefined visual patch of the target across frames.
\end{itemize}

\subsection{Region-Based Tracking}
Region-based tracking follows image patches with a \textbf{uniform appearance}. 
This method is computationally efficient and well-suited for real-time systems, providing a robust balance
between speed and discriminative power.

The standard implementation utilizes a \textbf{color histogram} as a visual signature. 
The workflow typically involves:
\begin{enumerate}
    \item Segmenting the region of interest (ROI) from the background.
    \item Computing the color distribution (histogram) of the ROI.
    \item Searching for the most similar histogram in the subsequent frame to locate the object.
\end{enumerate}

While effective for targets with distinctive colors, this approach is sensitive to illumination changes. 
To mitigate this, practitioners often employ the \textbf{HSV color space} (which separates luminance from chromaticity) 
or implement \textbf{dynamic model updates} to adapt the histogram over time.

\subsubsection{Part-Based Models}
A single global histogram may fail when tracking complex, non-rigid objects like humans. 
\textbf{Part-based models} decompose the object into several sub-regions (e.g., head, torso, and legs), 
each with its own local histogram.

This decomposition provides two major advantages:
\begin{itemize}
    \item \textbf{Robustness to Occlusion:} If the legs are hidden by an obstacle, 
    the system can maintain the track using the head and torso signatures.
    \item \textbf{Spatial Verification:} It prevents confusion between objects with similar colors but different 
    configurations. Take in consideration the case of two people, one wearing a white shirt and black pants, 
    and the other wearing a black shirt and white pants. A global histogram might struggle to differentiate them, 
    but a part-based model can leverage the spatial arrangement of colors to maintain accurate tracking.
\end{itemize}

\subsubsection{Implementation and Similarity Metrics}
The tracking process evaluates the similarity between the \textbf{tracked model} ($O^t$) 
in the current frame and a stored \textbf{reference model} ($O^r$). 

Common metrics for histogram comparison include:
\begin{itemize}
    \item \textbf{Histogram Intersection:} Measures the correlation between two distributions.
    \[ \cap (O^t, O^r) = \sum_{n=1}^{U} \min(O^t_{n}, O^r_{n}) \]
    This simply returns the area of the intersection of the two histograms, which is a measure of their similarity.
    \item \textbf{Sum of Squared Differences (SSD):} Measures the Euclidean distance between histogram bins.
    \[ SSD(O^t, O^r) = \sum_{n=1}^{U} (O^t_{n} - O^r_{n})^2 \]
    This instead returns the area of the difference between the two histograms, which is a measure of their dissimilarity.
\end{itemize}
\begin{comment}
    #TODO: insert image here of histogram comparison from the lesson
    1st method takes the minimum value of each corresponding bin and sums them up, giving as a result the area
    of the intersection of the two histograms. 
    2nd method takes the difference of each corresponding bin, squares it, and sums them up, giving as a result the area
    of the difference between the two histograms.
\end{comment}
To improve robustness, bins typically represent color ranges rather than individual 
values to account for sensor noise and minor lighting fluctuations.

\subsubsection{Addressing Environmental Noise: Shadows}
Shadows represent a significant challenge in region-based tracking because 
they alter the pixel values of the background, often leading to false positives in motion detection. 

To prevent a shadow from being tracked as a separate entity or distorting the object's histogram, 
we assume that a shadow primarily reduces \textbf{brightness} while leaving \textbf{chromaticity} (Hue and Saturation) relatively intact. 
By performing tracking in the HSV space and ignoring the Value ($V$) component, 
the system can become significantly more shadow-invariant, improving the robustness of the tracking process in outdoor 
environments or scenes with variable lighting conditions.

\subsubsection{Blob Extraction and Post-Processing}
Following motion detection, the resulting binary masks often contain "noise"—fragmented regions, 
isolated pixels, or holes within moving objects caused by imperfect segmentation. 
To transform these raw pixels into actionable data, we perform \textbf{blob extraction} (also known as Connected Component Labeling).

This process aggregates neighboring foreground pixels into coherent, 
labeled regions called \textbf{blobs}, which represent the physical objects in the scene. 
To ensure these blobs accurately reflect the geometry of the targets, the extraction process typically involves the following refinement steps:
\begin{itemize}
    \item \textbf{Morphological Filtering:} Applying \textit{erosion} to remove isolated noise (small pixel clusters) and \textit{dilation} to bridge gaps and fill internal holes within an object.
    \item \textbf{Connectivity Analysis:} Grouping pixels based on 4-way or 8-way adjacency to assign a unique ID to each spatially distinct moving entity.
    \item \textbf{Feature Extraction:} Once a blob is defined, we calculate its geometric properties, such as the \textbf{centroid}, \textbf{bounding box}, \textbf{area}, and \textbf{aspect ratio}.
\end{itemize}
By filtering out blobs that are too small to be objects of interest, we provide the downstream tracking algorithms with a clean, high-level representation of the scene.

\subsubsection{Target Association}
Once we have the blobs representing the detected objects, we need to associate them across consecutive frames to maintain the identity of each object over time.
This process is known as \textbf{target association} and is actually common to all tracking methods, not just region-based ones.

In general, detection is not carried out on a frame-by-frame basis, but rather on a \textbf{temporal window} of frames, since it can
be demanding in terms of computational resources. 
This means that the tracking algorithm must be able to handle cases where an object is not detected in every single frame, which adds complexity to the association process.

Correspondence is typically achieved through \textbf{data association algorithms} that match the detected blobs based on their spatial proximity and appearance features.
Such methods rely on the \textbf{optical flow} assumption, which states that the motion of objects between frames is small enough to allow for reliable matching based on their positions and appearances. 
It's important to keep in mind that cameras and frame rates must be chosen to satisfy this assumption, as large displacements can lead to tracking failures.

Different approaches to target association include:
\begin{itemize}
    \item \textbf{Nearest Neighbor:} Assigning each detected blob to the closest track based on spatial distance between centroids;
    \item \textbf{Overlapping:} Matching blobs based on the degree of overlap between their bounding boxes. 
    This has issues with scale changes;
    \item \textbf{Bounding Box - Overlap:} Similar to the previous method, but consider the bounding box of the track instead of the blob.
\end{itemize}
There are no best practices for target association, and the choice of method depends on the specific application and the characteristics of the scene being analyzed.

\subsubsection{Splitting and Merging}
Very similar but distinct problems that can arise during tracking are \textbf{splitting} and \textbf{merging}.

\paragraph{Splitting}
When multiple blobs are present in the same frame, we might encounter one of two problematic scenarios:
\begin{enumerate}
    \item \textbf{Single object - multiple blobs:} this occurs when an object is fragmented into several blobs due to imperfect segmentation.
    \item \textbf{Multiple objects - single blob:} this happens when two or more objects are merged into a single blob, often due to occlusion or proximity.
\end{enumerate}
These 2 situations are different faces of the same problem, which is called \textbf{splitting}. 
\begin{comment}
    #TODO: insert image here for example 1 and 2 from slides
\end{comment}
Before understanding which of the two cases we are facing, evidence need to be collected over a temporal window of frames.
This way, we can analyze the behavior of the blobs over time to determine whether they represent a single object or multiple objects.

\paragraph{Merging}
On the other hand, \textbf{merging} occurs when multiple objects are detected as a single blob, due to proximity or occlusion.
An example of this is when a person which is entering a scene. Initially, feet and arms might be detected as separate blobs.
As the person moves further into the scene, these blobs might merge into a single blob representing the whole body.
\begin{comment}
    #TODO: insert image here for example of merging from slides
\end{comment}

In both cases, its required to analyze the temporal evolution of the blobs to determine whether they represent a single object or multiple 
objects, and to adjust the tracking accordingly.

\subsubsection{Occlusion}
Moving object can be partially or fully occluded by each other when they are in close proximity.
This can lead to tracking failures, as one object disappears from the scene, and bigger blobs may appear, which
might have no resemblance to the original objects.

This is something that need to be handled before the occlusion occurs, by analyzing the proximity 
of the objects and predicting potential occlusions. 

Usually, the system suspends the tracking of the involved objects during the occlusion, and then re-identifies
them once they reappear as separate entities in the scene.

\subsection{Feature Based Tracking}
The idea behind feature-based tracking is to replace the information carried by the histogram of the whole object 
with a set of distinctive local \textbf{features}, extracted from the object itself.
These features are typically points of interest, such as corners, edges, or blobs, that can be reliably detected and tracked across frames.

Considering a set of images $A = \{I_1, I_2, ..., I_n\}$, we can extract a set of features $F_i$ from each image $I_i$.
The position of each feature in a given image can be represented as a vector of coordinates $f_i(x_i, y_i)$.
The objective of feature-based tracking is to determine the displacement vector $d_i = (dx_i, dy_i)$ that best estimates the movement of
feature $f_i$ from one frame to the next, such that:
\[ f_{i+1}(x_i + dx_i, y_i + dy_i) \]

\subsubsection{Good Features to Track}
Not all features are equally suitable for tracking. 
The \textbf{Good Features to Track} algorithm identifies features that are likely to be stable and distinctive across frames, making them ideal candidates for tracking.
This method is based on the concept of \textbf{corner detection}, which identifies points in the image where the intensity changes significantly in multiple directions.

For each pixel in the image $p(x, y)$ we compute the \textbf{structure tensor} $Z$:
\[ Z = \begin{bmatrix}
    \sum_{W} J_x^2 & \sum_{W} J_x J_y \\
    \sum_{W} J_x J_y & \sum_{W} J_y^2
\end{bmatrix} \]
where $J_x$ and $J_y$ are the gradients of the image in the x and y directions, respectively, 
and the summation is performed over a window $W$ centered around the pixel $p$.

The eigenvalues of the structure tensor, $\lambda_1$ and $\lambda_2$, provide a measure of the local intensity variation around the pixel.
A pixel is considered a good feature to track if the smallest eigenvalue $\lambda_{min} = \min(\lambda_1, \lambda_2)$ is above a certain threshold, 
indicating that changes are in multiple directions, which is characteristic of corners.

\subsubsection{Tracking Features}
Once we have identified the good features to track, we need to track their movement across frames.
Ideally, we would expect that the position of a feature in the next frame can be predicted by adding a displacement vector to its current position:
\[ I_i(Df - d) = I_{i+1}(f) \]
where $I_i, I_{i+1}$ are successive frames, $f$ is the position of the feature in the current frame, $D$ is the deformation matrix 
and $d$ is the displacement vector we want to estimate.

However, due to the noise, the equality does not hold perfectly. Assuming that the motion between frames is small, 
we can use a translational model to approximate the movement of the feature, which simplifies the problem to 
finding the displacement vector $d$ that minimizes the difference between the two images:
\[ \epsilon = \int \int_{W} [I_i(f - d) -  I_{i+1}(f)]^2 \omega(f) \; df \]
The value of $d$ that minimizes this error can be found by simply taking the derivative of $\epsilon$ with respect to $d$ and setting it to zero.